% !TeX root = ../main.tex
\section{结论}
\label{sec:conclusion}

本文围绕当前与未来混合精度负载的算力配比选择，提出组合式多精度点积微架构 DP-Puzzle，构建九种共享后端的配比设计。12\,nm、1\,GHz 下，$1:2:8$ 相比 $1:2:4$ 将 MXFP4 吞吐翻倍，面积增加 33.60\%，面效与能效分别提高 49.70\% 和 27.27\%，体现了精度专用增配与通路共享的收益。

联合面积、功耗与精度负载迁移的探索表明，高功耗模式会通过压缩部署规模影响其他精度的收益。在基准 TensorCore 面积占比 30\%、相同面积与 DP 池功耗预算下，$1:2:8$ 比 $1:4:8$ 提供高 60\% 的整芯片 FP4 算力，并在 DeepSeek-V4.1-Flash 的 1K--1M 上下文与六场景等权评测中取得最低平均计算时间，相对基线的加速比为 1.147 倍。该理想计算结果说明，保留部署密度与高精度算力的适度增配，能够适应多种低精度迁移路径；固定芯片的配比应由实现成本与跨负载收益共同决定。
